\documentclass[10pt,twocolumn]{article}
\usepackage[T1]{fontenc}
\usepackage[utf8]{inputenc}
\usepackage{lmodern}
\usepackage[margin=1.8cm,top=2.4cm,bottom=2cm,columnsep=0.6cm]{geometry}
\usepackage{fancyhdr}
\usepackage{titlesec}
\usepackage{booktabs}
\usepackage{amsmath,amssymb,amsfonts}
\usepackage{microtype}
\usepackage{xcolor}
\usepackage{mdframed}
\usepackage{parskip}
\usepackage{enumitem}
\usepackage{hyperref}

\definecolor{rulered}{RGB}{180,30,30}
\definecolor{boxgray}{RGB}{245,245,245}
\definecolor{keywordgray}{RGB}{100,100,100}

\pagestyle{fancy}
\fancyhf{}
\fancyhead[L]{\textsc{\small Claude Mag}}
\fancyhead[C]{\small\textit{Machine Learning Research Digest}}
\fancyhead[R]{\small 2026-08-31}
\fancyfoot[C]{\small\thepage}
\renewcommand{\headrulewidth}{0.8pt}

\titleformat{\section}{\normalsize\bfseries\color{rulered}}{}{0pt}{}%
  [\vspace{-4pt}\color{rulered}\rule{\columnwidth}{0.4pt}]
\titlespacing{\section}{0pt}{8pt}{4pt}

\hypersetup{colorlinks=true,urlcolor=rulered,linkcolor=black}

\begin{document}
\setlength{\parindent}{0pt}
\setlength{\parskip}{4pt}

{\small\color{keywordgray}\textbf{Keywords:}
  mechanistic interpretability \textbullet{}
  causal circuits \textbullet{}
  post-training \textbullet{}
  LLM fine-tuning}\par\vspace{4pt}

{\LARGE\bfseries\raggedright
  Circuit Condensation: Post-Training that
  Concentrates a Behavior's Causal Circuit}\par\vspace{4pt}

{\large\itshape\raggedright
  SFT Universally Condenses Causal Circuits While RL
  Expands Multi-Step Reasoning Circuits}\par\vspace{6pt}

{\small\textbf{By} Sai Adith Senthil Kumar}\par\vspace{2pt}

{\small\color{keywordgray}arXiv:2608.27254 \textbullet{} Preprint, August 2026}
\par\vspace{2pt}\noindent{\color{rulered}\rule{\columnwidth}{0.6pt}}\vspace{6pt}

Does post-training reorganise the internal circuits of a language
model in interpretability-friendly ways? Circuit Condensation answers
this question empirically across eight behaviours and three post-training
regimes (SFT, GRPO, and base). Using ACDC and edge attribution patching,
the paper shows that SFT consistently shrinks and focuses causal circuits
(+20\% faithfulness, $-$31\% component count), while RL post-training
condenses simple circuits but expands complex multi-step ones — a
divergence with direct implications for scalable interpretability.

\begin{mdframed}[backgroundcolor=boxgray,linecolor=rulered,linewidth=1pt,
    innerleftmargin=6pt,innerrightmargin=6pt,
    innertopmargin=6pt,innerbottommargin=6pt]
\textbf{\small AT A GLANCE}\par\vspace{4pt}
\begin{description}[leftmargin=0pt,labelsep=4pt,itemsep=2pt,parsep=0pt,
    font=\small\bfseries]
  \item[Problem:] It is unknown whether post-training concentrates or
    disperses the causal circuits responsible for specific behaviors,
    with direct consequences for scalable mechanistic interpretability.
  \item[Why it matters:] If post-training concentrates circuits,
    automated circuit analysis becomes easier; if it disperses them,
    interpretability may degrade even as capability improves.
  \item[Method:] ACDC and Edge Attribution Patching applied before and
    after SFT, GRPO, and combined post-training across 8 representative
    LLM behaviors.
  \item[Result:] SFT: $-$31\% circuit size, $+$0.15 edge faithfulness.
    GRPO: condenses factual circuits but expands multi-step reasoning
    circuits by $+$47\%.
  \item[Evidence:] Effect is robust across learning rates and dataset
    sizes for SFT; behavior-type (simple vs.\ multi-step) determines
    RL's direction of effect.
\end{description}
\end{mdframed}

\section{The Big Picture}

Mechanistic interpretability aims to identify the sparse subgraph
of a neural network's computation that causally produces a specific
behaviour — its \emph{circuit}. Progress in circuit discovery has
revealed that transformers implement interpretable algorithms in
relatively compact circuits: the indirect object identification
circuit in GPT-2 involves $\sim$26 attention heads; the greater-than
circuit involves $\sim$10.

But circuits are always identified in \emph{base} models. Post-training
dramatically changes model behaviour. A natural question follows:
does post-training make circuits more concentrated (easier to identify
and edit) or more dispersed (harder)? This question is not merely
academic — interpretability tooling built for base models may
perform very differently on fine-tuned models used in deployment.

Circuit Condensation provides the first systematic answer.

\section{Why Existing Approaches Fall Short}

\textbf{Static analysis:} Prior mechanistic interpretability work
analyses circuits in fixed (base or RLHF'd) models but does not
compare circuit structure across training stages. There is no prior
data on whether post-training increases or decreases circuit sparsity.

\textbf{Limited behaviour coverage:} Some work examines single
behaviours (e.g., factual recall) before and after fine-tuning,
but the generality of any observed effect is unclear. Circuit
Condensation covers eight behaviours of varying complexity to
identify patterns.

\textbf{Confounded by capability:} If a model becomes more accurate
at a behaviour after fine-tuning, circuit quality may change even if
circuit structure is unchanged. The paper controls for capability
by measuring circuit faithfulness at matched accuracy levels.

\section{Core Mechanism}

\textbf{Circuit discovery methods:}

\textbf{ACDC (Automated Circuit Discovery)} iteratively ablates
edges in the computation graph — replacing activations with
corrupted versions from a distribution matched to the task —
and retains an edge only if its ablation decreases performance
by more than threshold $\epsilon$. The result is a sparse
\emph{minimal circuit} sufficient to perform the behaviour.

\textbf{Edge Attribution Patching (EAP)} scores each edge by
a first-order approximation:
\[
  \text{EAP}(e) = \frac{\partial \mathcal{L}}{\partial h_e}
    \cdot (h_e^\text{clean} - h_e^\text{corrupt})
\]
where $h_e$ is the activation at edge $e$, and $\mathcal{L}$ is
the task loss. EAP enables fast ranking of all edges without
exhaustive ablation.

The paper uses ACDC to identify circuits (trusted, slower) and
EAP to validate patterns (fast, scalable to larger models).

\textbf{Circuit metrics:}
\begin{itemize}[itemsep=2pt,parsep=0pt]
  \item \textbf{Circuit size}: number of attention heads and MLP
    layers with non-zero ACDC edges
  \item \textbf{Edge faithfulness}: fraction of task performance
    recovered when running only the circuit (with other components
    ablated)
  \item \textbf{Modularity}: intersection-over-union of circuit
    components across behaviors (lower = more behavior-specific)
\end{itemize}

\section{Technical Formulation}

Let $C_b(\theta)$ be the ACDC circuit for behavior $b$ under
parameters $\theta$. The paper tracks three metrics across $\theta$:
\[
  \text{Size}(C_b) = |V(C_b)| + |E(C_b)|
\]
\[
  \text{Faith}(C_b) = \frac{\text{Perf}_{C_b}}{\text{Perf}_\theta}
\]
\[
  \text{Mod}(b_1, b_2) = 1 - \frac{|C_{b_1} \cap C_{b_2}|}{|C_{b_1} \cup C_{b_2}|}
\]
where $\text{Perf}_{C}$ is task performance with only circuit $C$
active. Comparisons are made between $\theta_\text{base}$,
$\theta_\text{SFT}$, and $\theta_\text{RL}$.

\section{Learning or Inference Procedure}

For each of 8 behaviors, the paper trains three model versions:
\begin{enumerate}[itemsep=2pt,parsep=0pt]
  \item Base model (no post-training)
  \item SFT on 500 behavior-specific demonstrations
  \item GRPO on the same 500 problems with scalar correctness rewards
\end{enumerate}

Circuit discovery is run at each checkpoint using ACDC (threshold
$\epsilon = 0.01$) on 100 held-out task instances. EAP is run on
500 instances to validate ACDC findings.

\section{What the Guarantee Says}

No formal convergence guarantees are given. The empirical claim
is that the observed patterns (SFT condenses; RL behavior depends
on task complexity) are robust across model sizes (1B--7B parameters
tested), learning rates, and dataset sizes.

\section{Experimental Findings}

\begin{itemize}[itemsep=2pt,parsep=0pt]
  \item \textbf{SFT (all behaviors):} Average circuit size decreases
    by $-$31\%; edge faithfulness increases from 0.74 to 0.89
  \item \textbf{GRPO (simple behaviors: factual recall, IOI):}
    Circuit size decreases $-$28\%; faithfulness $+$0.11 — similar
    to SFT
  \item \textbf{GRPO (multi-step behaviors: arithmetic, logical):}
    Circuit size \emph{increases} $+$47\%; faithfulness \emph{decreases}
    $-$0.06 — opposite direction from SFT
  \item \textbf{Modularity:} SFT increases modularity (circuits become
    more behavior-specific); GRPO has mixed effects
  \item \textbf{Combined SFT+GRPO:} SFT condensation is partially
    preserved after subsequent GRPO for simple behaviors, but
    overridden for multi-step behaviors
\end{itemize}

\section{Ablations and Interpretation}

\textbf{Learning rate:} SFT condensation is robust across learning
rates from $10^{-5}$ to $10^{-3}$; the magnitude of size reduction
varies slightly but direction is consistent.

\textbf{Dataset size:} Condensation occurs even with as few as 50
demonstrations; more demonstrations yield larger faithfulness gains.

\textbf{Why does RL expand multi-step circuits?} The paper hypothesises
that RL via GRPO improves multi-step performance by recruiting
additional attention heads to handle credit propagation over long
chains. This is consistent with findings from other work showing
that RL fine-tuning on reasoning tasks induces new attention head
specialisations.

\textbf{Implications for circuit-based editing:} For SFT-trained
models, targeted activation patching and component-level editing
become more reliable because circuits are more concentrated and
faithful. For RL-trained models used on complex reasoning tasks,
the expanded circuits make editing harder — a previously unexplored
cost of RL post-training.

\section*{Reference}

Sai Adith Senthil Kumar.
\textbf{Circuit Condensation: Post-Training that Concentrates
a Behavior's Causal Circuit.}
arXiv:2608.27254, August 2026.
\url{https://arxiv.org/abs/2608.27254}

\end{document}
